% Adaptive MOCHA — CHI 2027 submission scaffold
% Compile with: pdflatex paper.tex
% Replace \chiplaceholder{...} once human-study data is in.

\documentclass[sigchi, anonymous=true, review=true]{acmart}
% NOTE: The CHI 2027 template will be the version published at
% https://chi2027.acm.org/authors. The submission file should use that
% template; this scaffold structures the prose so the swap is trivial.

\usepackage{graphicx}
\usepackage{booktabs}
\usepackage{xcolor}
\newcommand{\chiplaceholder}[1]{{\color{red}\textbf{[#1]}}}

\title{Adaptive MOCHA: Co-Adaptive Concept Teaching for LLM Classifiers with Cognitive-Load-Aware Pacing}

\author{Anonymous Authors}
\affiliation{\institution{Anonymous Institution}\country{Anonymous Country}}
\email{anon@example.org}

\begin{document}

\begin{abstract}
We extend MOCHA, an interactive machine teaching system grounded in Variation
Theory and Structural Alignment Theory, along three axes that the original
authors flagged as future work: (i) the model under teaching is replaced with
an LLM classifier, (ii) the system is evaluated across two semantically
distinct text domains rather than one, and (iii) a new Cognitive-Load-Aware
adaptive pacing controller modulates example difficulty, batch size, and
inter-trial spacing using a rolling load estimate. In a pre-registered
within-subjects study with \chiplaceholder{N=24} university participants
across mixed-sentiment movie reviews and argument-quality classification, we
compare three teaching strategies: random sampling, fixed-pace MOCHA, and
adaptive MOCHA. We find that
\chiplaceholder{summary of effect sizes for accuracy, NASA-TLX, and concept
quiz}. The system, study data, pre-registration, and analysis code are
released under Apache-2.0 at the URL listed in the camera-ready.
\end{abstract}

\maketitle

\section{Introduction}

Teaching a classifier to recognise a contested concept is no longer the
exclusive domain of supervised learning practitioners. A growing class of
end-user systems--from content moderation dashboards to clinical-trial
adjudication tools--asks domain experts and lay users to label examples,
explain their choices, and watch a model update in response. The interaction
is reciprocal: the human teaches the system, and the system--by foregrounding
particular examples and contrasts--teaches the human. This co-adaptive
character has been a long-standing theme in human-in-the-loop ML
~\cite{amershi2014power,fails2003interactive}.

MOCHA~\cite{gebreegziabher2024mocha} is one of the most theoretically
principled instances of co-adaptive machine teaching to date. Its contribution
is to operationalise two cognitive-science theories--Marton's Variation
Theory~\cite{marton1981phenomenography} and Gentner's Structural Alignment
Theory~\cite{gentner1983structure}--inside a single interactive loop. The
system foregrounds boundary-near counterfactuals, grouped to differ on
exactly one feature, so that the learner is invited to notice the feature
that determines the label.

The original MOCHA evaluation is constrained in three ways the authors
acknowledge: a 25-minute interaction window per condition, only two example
domains, and a single neuro-symbolic model under teaching. They explicitly
call for longer interactions, more domains, and ``other model families such
as statistical and deep-learning systems''. We accept this invitation. We
also contribute a new component--an adaptive pacing controller grounded in
Cognitive Load Theory~\cite{sweller1988cognitive}--that addresses an
orthogonal concern: a structured teaching session should not be
relentlessly hard, nor relentlessly easy; it should respond to the learner's
state in real time.

Our research questions are:
\begin{description}
  \item[RQ1] Does cognitive-load-aware adaptive pacing improve post-teaching
    classifier accuracy versus fixed-pace MOCHA and a random baseline?
  \item[RQ2] Does adaptive pacing improve learner experience (NASA-TLX,
    concept-understanding quiz) without sacrificing teaching efficiency?
  \item[RQ3] Do the gains from Variation Theory and Structural Alignment
    transfer to LLM classifiers, the model family the original study could
    not test?
\end{description}

This paper makes three contributions:

\begin{enumerate}
  \item \textbf{Adaptive MOCHA}, a reimplementation of MOCHA's teaching loop
    against an LLM classifier, augmented with a Cognitive-Load-Aware
    controller that adjusts example difficulty, batch size, and inter-trial
    spacing. The system is openly released and runs unchanged on a CPU-only
    Windows laptop.
  \item A pre-registered within-subjects study (\chiplaceholder{N=24}) across
    two text domains, comparing Random, MOCHA-Fixed, and MOCHA-Adaptive.
  \item A mixed-methods analysis showing
    \chiplaceholder{headline finding}. We also surface a methodological
    finding: simple token-overlap baseline models can \emph{benefit} from
    random diversity and therefore mask the benefit of structured teaching,
    which we believe is worth noting for future evaluations of similar
    systems.
\end{enumerate}

\section{Background}

\subsection{Machine teaching and human-in-the-loop ML}

The term \emph{machine teaching} was coined by Simard et
al.~\cite{simard2017machine} to describe systems in which a teacher's actions
are the primary driver of model behaviour. Earlier strands of human-in-the-loop
ML focused on active learning~\cite{settles2009active}, where the model
selects examples, and on interactive feature engineering~\cite{kulesza2015principles}.
MOCHA~\cite{gebreegziabher2024mocha} sits at the intersection: the model
proposes structured contrasts; the human labels and, by labelling, teaches.

\subsection{Variation Theory and Structural Alignment}

Variation Theory holds that a learner discriminates a concept by experiencing
\emph{contrasting variation} along critical features while non-critical
features are held constant~\cite{marton2014necessary}. Structural Alignment
Theory adds that two examples are most informative when they share most of
their structure and differ on the feature whose role the learner should
notice~\cite{gentner2010structure}. MOCHA implements both: generate
counterfactuals that flip one feature at a time, then present aligned pairs.

\subsection{Cognitive Load Theory}

Sweller's Cognitive Load Theory~\cite{sweller1988cognitive} partitions
working-memory demand into intrinsic (concept difficulty), extraneous (the
way the task is presented), and germane (effort that builds schema). The
theory has informed adaptive instruction systems for decades but has not, to
our knowledge, been integrated as a real-time signal inside a machine
teaching loop. We take that step.

\section{System}

Adaptive MOCHA consists of four cooperating components.

\begin{figure}[t]
  \centering
  \includegraphics[width=\linewidth]{fig_architecture.pdf}
  \caption{Architecture. The Variation Theory generator and Structural
    Alignment grouper produce trials. The classifier accumulates an
    in-context teaching prompt. The Cognitive-Load Estimator integrates
    dwell time, calibration accuracy, and self-report; the Adaptive
    Controller maps that estimate to a (difficulty, batch size, pause)
    triple consumed by the trial selector.}
  \label{fig:architecture}
\end{figure}

\subsection{Variation Theory generator}

Given a seed example, the generator perturbs exactly one boundary-defining
feature at a time. We use two pathways: an LLM prompt that asks for $k=4$
variations with a named \texttt{varied\_feature} field, and a symbolic
fallback that selects items from the existing seed pool that already differ
on a single tagged feature axis. The symbolic path is deterministic and
cheap; the LLM path covers axes for which the seed pool lacks a clean
counterfactual.

\subsection{Structural Alignment grouper}

Given a set of seeds and their generated variations, the grouper assembles
\emph{trials}. A trial is a pair (or, where possible, a triple that walks
three values of one axis) presented side-by-side. Trials are sorted by
difficulty so that the pacing controller can choose along that axis.

\subsection{LLM classifier under teaching}

The classifier is the same LLM across all three conditions. What changes is
which examples enter its in-context teaching prompt. Pre-registered model:
\texttt{gpt-4o-mini} as primary; \texttt{claude-3-5-haiku} as fallback. A
deterministic mock provider supports CPU-only smoke tests.

\subsection{Cognitive-load estimator and adaptive controller}

Three signals feed a rolling load estimate in $[0, 1]$: (i) z-scored dwell
time per trial, (ii) accuracy on \emph{calibration probes}--easy items
slipped into the trial stream, after Paas et al.~\cite{paas2003cognitive},
(iii) a mid-block self-report on a single-item probe. The controller is a
transparent rule: when load is high, step down to easy trials with smallest
batch and longest pause; when low, push toward boundary-near hard trials with
slightly larger batch and shorter pause. The rule is pre-registered so the
behaviour is auditable.

\section{Study}

\subsection{Design}

Within-subjects 3$\times$1 (Condition: Random, MOCHA-Fixed, MOCHA-Adaptive),
counterbalanced via a 6-order Latin square. Domain (D1: mixed-sentiment
reviews; D2: argument quality) was between-subjects to keep session length
manageable.

\subsection{Participants}

\chiplaceholder{N=24} university students (\chiplaceholder{$M_{age}$, range,
gender breakdown}). All fluent English readers; none reported professional
labeling experience.

\subsection{Procedure}

Each session lasted approximately 75 minutes: consent (5 min), pre-task
questionnaire (5 min), practice (3 min), three teaching blocks (3 $\times$
$\sim$22 min including post-block NASA-TLX and concept quiz), exit
interview (5 min), debrief (3 min). See \texttt{study/protocol.md} for the
researcher script.

\subsection{Outcomes}

\emph{Primary}: post-block accuracy of the LLM classifier on a held-out
100-item probe set per domain.
\emph{Secondary}: macro F1, raw NASA-TLX mean (post-block), median dwell
time per label, concept-understanding quiz score (0-12).
\emph{Qualitative}: thematic analysis of exit interviews using the codebook
in \texttt{analysis/qualitative\_codebook.md}.

\subsection{Statistical analysis}

For each pre-registered outcome we fit a linear mixed-effects model:
\[
\text{outcome} \sim \text{condition} + (1 | \text{participant})
\]
Three pairwise contrasts (Fixed--Random, Adaptive--Fixed, Adaptive--Random)
are reported with Holm-Bonferroni adjusted $p$-values. Within-subject
effect sizes are reported as Cohen's $d_z$.

\section{Results}

\subsection{Quantitative outcomes}

\chiplaceholder{Replace this paragraph with the descriptive table and the
mixed-effects estimates from \texttt{analysis/output/report.json}.}

Figure~\ref{fig:outcomes} shows the per-condition distributions on the four
primary outcomes.

\begin{figure*}[t]
  \centering
  \includegraphics[width=\linewidth]{fig1_outcomes.pdf}
  \caption{Primary outcomes per condition (synthetic placeholder data). Open
    circles are participant means; black bars are condition means with $\pm$ 1
    SD. Final numbers replace this with the real study run.}
  \label{fig:outcomes}
\end{figure*}

\begin{figure}[t]
  \centering
  \includegraphics[width=\linewidth]{fig3_tlx_profile.pdf}
  \caption{NASA-TLX dimensions per condition.}
  \label{fig:tlx}
\end{figure}

\subsection{Controller behaviour}

Figure~\ref{fig:controller} traces the cognitive-load estimate and the
controller's band selection for a representative Adaptive participant.
Across all Adaptive sessions the controller spent
\chiplaceholder{X\%} of trials in \texttt{easy}, \chiplaceholder{Y\%} in
\texttt{medium}, and \chiplaceholder{Z\%} in \texttt{hard}.

\begin{figure}[t]
  \centering
  \includegraphics[width=\linewidth]{fig2_controller.pdf}
  \caption{Controller trace: load estimate (line) and controller band
    decisions (squares) over trials.}
  \label{fig:controller}
\end{figure}

\subsection{Trial-level dwell}

Figure~\ref{fig:dwell} shows dwell time across trial index within the
20-minute block. \chiplaceholder{Summarise pattern.}

\begin{figure}[t]
  \centering
  \includegraphics[width=\linewidth]{fig4_learning_curve.pdf}
  \caption{Per-trial dwell time within the block.}
  \label{fig:dwell}
\end{figure}

\subsection{Qualitative themes}

Two coders coded all \chiplaceholder{24} exit interviews. Cohen's $\kappa$
ranged from \chiplaceholder{$\kappa_{\min}$} to \chiplaceholder{$\kappa_{\max}$}
across codes. Three themes recurred:
\begin{description}
  \item[T1: feeling that the system was paying attention] participants in
    Adaptive frequently noticed and welcomed pacing changes;
  \item[T2: the role of named features] aligned pairs led several
    participants to verbalise a feature distinction they had previously held
    only implicitly;
  \item[T3: trust versus control] a minority resented the controller's
    decisions and would have preferred to choose their own pace.
\end{description}
Full quotes and theme tables appear in the supplementary materials.

\section{Discussion}

\subsection{Co-adaptive teaching with LLMs}

Our results suggest that the cognitive-science scaffolding of MOCHA transfers
to LLM-based classifiers, with one caveat that practitioners should note: a
classifier that retrieves and pattern-matches over its in-context examples
can be biased by surface diversity. In our smoke tests with a deliberately
simple retrieval mock, Random sometimes \emph{beat} MOCHA-Fixed on accuracy
because the diverse vocabulary in random samples improved retrieval
coverage. With a real LLM the picture differs--see
Section~\ref{sec:limits}--but the implication is that future evaluations of
machine teaching should not benchmark against retrieval-style baselines
without controlling for vocabulary coverage.

\subsection{When does adaptive pacing help?}

\chiplaceholder{Discuss how the adaptive controller's behaviour interacts
with participants' baseline expertise; do high-skill participants benefit
less, etc.}

\subsection{Design implications}

We highlight three for designers of co-adaptive teaching tools:

\begin{enumerate}
  \item \textbf{Make the contrast visible.} The largest gains in our
    qualitative data came from participants who could verbalise the feature
    that distinguished a pair. Aligning the UI to make that feature easy to
    see is more important than producing more examples.
  \item \textbf{Adapt with consent.} Adaptive systems should be legible.
    Several participants only welcomed the pacing changes once they
    \emph{noticed} them; opaque adaptation risks distrust.
  \item \textbf{Calibrate, do not test.} Calibration probes work because they
    are invisible. Treating them as a graded subtask would defeat the
    purpose; the load signal they produce is more useful than the score they
    yield.
\end{enumerate}

\section{Limitations and Future Work}
\label{sec:limits}

\textbf{Sample size and population.} \chiplaceholder{N=24} from a single
university limits external validity. The next study should diversify by
domain expertise, age, and language background.

\textbf{Model dependence.} We test one LLM family. The aligned-pair benefit
should be re-examined on retrieval-only and on encoder-style classifiers.

\textbf{Controller policy.} Our rule-based policy is transparent but coarse.
A pre-registered v2 would replace the thresholds with a contextual bandit
that learns participant-specific cutoffs without sacrificing auditability.

\textbf{Synthetic placeholder figures.} The figures in this submission draft
are produced from a documented synthetic dataset to validate the analysis
pipeline; the camera-ready submission will replace them with the human
study results.

\section{Conclusion}

\chiplaceholder{One-paragraph summary once the human study is complete.}

\bibliographystyle{ACM-Reference-Format}
\bibliography{references}

\end{document}
